\section{Previous Work}
\label{sec:hao-background}

Direct-P starts from HAO AI Lab's FP4 FA4 implementation
\cite{hao2026fp4flashattention}.  HAO already provides a two-query pipeline
that assigns data movement, matrix products, softmax, and output correction to
specialized warps.  We keep that outer schedule.  Our forward contribution is
the narrower stage that turns a completed score fragment into the scaled
probability operand consumed by the value product.  This section first
explains the tiled algorithm, then states exactly which scheduling and storage
decisions we inherit.

\subsection{Why FlashAttention uses an online softmax}

For one query row, exact attention is
\begin{equation}
  O_i=\frac{\sum_j e^{z_{ij}}V_j}{\sum_j e^{z_{ij}}},
  \qquad z_{ij}=Q_iK_j^\mathsf{T}/\sqrt d.
  \label{eq:exact-attention}
\end{equation}
FlashAttention visits the keys in tiles.  After a tile with maximum
$\widetilde m_i$ arrives, it updates a running maximum $m_i$, denominator
$\ell_i$, and output $o_i$:
\begin{align}
  m_i' &= \max(m_i,\widetilde m_i), \\
  \ell_i' &= e^{m_i-m_i'}\ell_i+
    \sum_{j\in B}e^{z_{ij}-m_i'}, \\
  o_i' &= e^{m_i-m_i'}o_i+
    \sum_{j\in B}e^{z_{ij}-m_i'}V_j.
  \label{eq:online-recurrence}
\end{align}
This recurrence avoids storing the full score and probability matrices in
HBM.  Its cost is a strict dependency: a score tile must be reduced and turned
into probabilities before those probabilities can be multiplied by $V$.

On data-center Blackwell GPUs, asynchronous \tcgen{} operations accumulate
32-bit floating-point (FP32) score and output tiles in tensor memory (TMEM).
The Tensor Memory Accelerator (TMA) supplies shared-memory operands.  Separate
warps issue matrix operations, compute softmax, and correct the online output
\cite{nvidia2026ptx,zadouri2026flashattention4}.

\subsection{Implementation baseline}

HAO's implementation supports block-scaled FP4 QK with BF16 or FP8 PV, as
well as a stabilized path that uses NVFP4 throughout attention.  Its
implementation and Attn-QAT, an attention quantization-aware-training study,
identify online P quantization and scale movement as central costs
\cite{hao2026fp4flashattention,zhang2026attnqat}.

SageAttention3 contributes related numerical ideas, including Q/K smoothing,
two-level P scaling, and reuse of the block maximum during online softmax
\cite{zhang2025sageattention3}.  It targets a consumer Blackwell interface
that differs from the asynchronous tensor-core and tensor-memory interface on
the data-center Blackwell GPUs evaluated here \cite{nvidia2026ptx}.  We use
those numerical ideas as context, but inherit the execution schedule from HAO
and FA4.

\subsection{Two-query execution model}

For our purposes, HAO's key contribution is a complete ownership plan for the
Blackwell pipeline \cite{hao2026fp4flashattention}.  At head dimension 128
(D128), one 16-warp cooperative thread array (CTA) advances two 128-row query
tiles (M128), called stage 0 and stage 1.  A load warp keeps K and V ahead in a
TMA-fed shared-memory ring.  One warp issues ordered QK and PV matrix
multiply--accumulate (MMA) operations.  Two softmax warpgroups (WGs), each
made of four warps, process the query stages in ping-pong fashion.  Separate
warps correct the online state and store the final output.

Small hardware barriers act as event counters rather than stopping the whole
CTA.  They announce ``score ready'', ``first P half ready'', ``tail P ready'',
and ``output released''.  Algorithm~\ref{alg:hao-pipeline} states the retained
ownership protocol.  Operations for $q=0$ and $q=1$ interleave whenever their
dependencies allow; the table gives the required order for one score bank,
not a serialized whole-CTA loop.

N32 denotes a 32-column score or probability fragment.  K$x$ denotes an
$x$-element matrix-product inner dimension, so K64 consumes two adjacent N32
fragments.

\begin{algorithm}[htbp]
\caption{Two-query score-to-PV ownership protocol inherited from HAO}
\label{alg:hao-pipeline}
\small
\begin{tabularx}{\textwidth}{@{}r p{.18\textwidth} X@{}}
\toprule
& Owner & Event for query stage $q\in\{0,1\}$ \\
\midrule
1 & Load warp & Prefetch the next K/V tile into the shared-memory ring. \\
2 & MMA warp & Wait until score bank $S_q$ is free; issue QK into $S_q$ and publish \emph{score ready}. \\
3 & Softmax WG $q$ & Read N32 quarters Q0 and Q1, create packed P and scales in retired score addresses, then publish the first legal K64 operand. \\
4 & MMA warp & Observe the first-half event and issue asynchronous $PV_{q,0}$ into permanent output bank $O_q$. \\
5 & Softmax WG $q$ & Process Q2 and Q3 and publish the tail K64 operand. \\
6 & MMA warp & Issue $PV_{q,1}$; meanwhile issue legal work for stage $1-q$. \\
7 & Correction WG & Update online state; the epilogue normalizes and stores when the key loop ends. \\
8 & MMA warp & After tail PV has consumed the overlay, publish $S_q$ as free for the next QK tile. \\
\bottomrule
\end{tabularx}
\end{algorithm}

The two-query topology matters at high head count because one CTA performs two
query jobs while reusing the same K/V stream.  A one-query topology creates
twice as many independent CTAs and can add a second scheduling wave after all
streaming multiprocessors are occupied.  A third query stage would provide
more look-ahead, but the TMEM layout leaves no bank for it.

\subsection{Why the D128 layout cannot look farther ahead}

SM100 and the tested SM103 path expose 512 logical TMEM columns. A D128 FP32
score tile or output accumulator uses 128 columns. Two score banks and two
output banks therefore consume the entire allocation.

\begin{figure}[htbp]
\centering
\begin{tikzpicture}[x=0.044\textwidth,y=0.72cm,font=\small]
  \draw[thick] (0,0) rectangle (20,1);
  \fill[scoreblue!22] (0,0) rectangle (5,1);
  \fill[scoreblue!36] (5,0) rectangle (10,1);
  \fill[outputgreen!26] (10,0) rectangle (15,1);
  \fill[outputgreen!40] (15,0) rectangle (20,1);
  \draw (5,0) -- (5,1);
  \draw (10,0) -- (10,1);
  \draw (15,0) -- (15,1);
  \node at (2.5,.55) {$S_0$: score stage 0};
  \node at (7.5,.55) {$S_1$: score stage 1};
  \node at (12.5,.55) {$O_0$: output stage 0};
  \node at (17.5,.55) {$O_1$: output stage 1};
  \node[below] at (0,0) {0};
  \node[below] at (5,0) {128};
  \node[below] at (10,0) {256};
  \node[below] at (15,0) {384};
  \node[below] at (20,0) {512};
  \node[align=center,text=inkgray,font=\scriptsize,text width=.40\textwidth]
    at (5,-1.05)
    {temporary: QK score, then packed $P$ and scale pages};
  \node[align=center,text=inkgray,font=\scriptsize,text width=.40\textwidth]
    at (15,-1.05)
    {long lived: FP32 online output accumulators};
\end{tikzpicture}
\caption{D128 TMEM layout inherited from HAO. P and its instruction-facing
scales reuse score storage only after that score fragment is retired. There
is no unowned 128-column bank.}
\label{fig:tmem-layout}
\end{figure}

HAO reuses each score bank as
\begin{equation}
  \text{QK score}\rightarrow\text{read N32 fragment}
  \rightarrow\text{P/scale overlay}\rightarrow\text{PV consume}
  \rightarrow\text{free}.
  \label{eq:score-overlay}
\end{equation}
Writing the next QK tile too early destroys P while PV still needs it.
Waiting leaves the matrix issuer idle.  The scheduling problem is therefore a
storage-ownership handoff, not merely a shortage of barrier signals.  A
barrier can report that storage is free; it cannot create another destination.

\subsection{Why probability tiles are published in halves}

QK still writes one 128-row by 128-column (M128$\times$N128) FP32 score tile.
``Quartering'' means
that a softmax warpgroup reads and transforms four N32 fragments,
\begin{equation}
  S=[Q0\mid Q1\mid Q2\mid Q3],\qquad
  Qk\in\mathbb R^{128\times32},
\end{equation}
not that QK becomes four smaller matrix products. The score allocation remains
128 columns.

The scaled-FP4 PV instruction used by the retained HAO path consumes K64
rather than K32. Two adjacent quarters must therefore be complete before
useful matrix work can issue:
\begin{equation}
  (Q0,Q1)\rightarrow PV_{K64}^{(0)},\qquad
  (Q2,Q3)\rightarrow PV_{K64}^{(1)}.
  \label{eq:k64-pairs}
\end{equation}
Quartering reduces the live register fragment and lets P overwrite score
addresses that are no longer needed.  It does not reduce the score bank's
TMEM allocation.  Its performance value is the event after Q1: first-half PV
can begin while Q2/Q3 and the other query stage provide independent work.

\begin{figure}[htbp]
\centering
\begin{tikzpicture}[
  x=1.18cm,y=.82cm,font=\small,
  box/.style={draw,rounded corners=1.5pt,minimum width=1.22cm,
              minimum height=.56cm,align=center},
  dep/.style={-{Stealth[length=2mm]},thick}]
  \node[box,fill=scoreblue!26] (s0) at (0,1.5) {QK 0};
  \node[box,fill=probred!17] (a0) at (1.5,1.5) {Q0};
  \node[box,fill=probred!23] (a1) at (2.8,1.5) {Q1};
  \node[box,fill=outputgreen!25,minimum width=1.55cm] (p0) at (4.25,1.5) {PV 0a};
  \node[box,fill=probred!28] (a2) at (5.75,1.5) {Q2};
  \node[box,fill=probred!34] (a3) at (7.05,1.5) {Q3};
  \node[box,fill=outputgreen!38,minimum width=1.55cm] (p1) at (8.5,1.5) {PV 0b};

  \node[box,fill=scoreblue!38] (s1) at (1.0,0) {QK 1};
  \node[box,fill=probred!17] (b0) at (2.5,0) {Q0};
  \node[box,fill=probred!23] (b1) at (3.8,0) {Q1};
  \node[box,fill=outputgreen!25,minimum width=1.55cm] (p2) at (5.25,0) {PV 1a};
  \node[box,fill=probred!28] (b2) at (6.75,0) {Q2};
  \node[box,fill=probred!34] (b3) at (8.05,0) {Q3};
  \node[box,fill=outputgreen!38,minimum width=1.55cm] (p3) at (9.5,0) {PV 1b};

  \draw[dep] (s0)--(a0); \draw[dep] (a0)--(a1);
  \draw[dep] (a1)--(p0); \draw[dep] (a1)--(a2);
  \draw[dep] (a2)--(a3); \draw[dep] (a3)--(p1);
  \draw[dep] (s1)--(b0); \draw[dep] (b0)--(b1);
  \draw[dep] (b1)--(p2); \draw[dep] (b1)--(b2);
  \draw[dep] (b2)--(b3); \draw[dep] (b3)--(p3);
\end{tikzpicture}
\caption{Dependency sketch for the two query stages. Box width is not
measured time. Each first PV half waits for two N32 producer quarters; work
from the other stage covers part of that delay.}
\label{fig:quarter-pipeline}
\end{figure}

We inherit this outer pipeline, tensor-memory lifecycle, and publication
protocol unchanged.  They also expose the remaining problem: PV cannot begin
until softmax has constructed, scaled, packed, and published two adjacent
probability fragments.  The next section separates this timing constraint
from the numerical range constraint.
